% GENERATED FILE. Edit canonical lessons, then run npm run generate:books.
% canonical-source-hash: fnv1a64:d0c0aebdf24c0828
% canonical-lessons: BN-C150-sondeh-kora, BN-C150-torko-kora, BN-C150-chaliye-jaoya, BN-C150-hacha, BN-C150-kasha

\chapter{To doubt, To argue, To carry on, To sneeze, To cough}
\label{ch:bn-to-doubt-to-argue-to-carry-on-to-sneeze-to-cough}
\hlchaptermodality{150}

\begin{chapteropening}
\textbf{By the end of this chapter:} \emph{I can say to doubt, to argue, to carry on, to sneeze and to cough.}

Complete the last lesson of chapter 150: I can say to doubt, to argue, to carry on, to sneeze and to cough.
\end{chapteropening}

\section[\texorpdfstring{\bn{সন্দেহ} \bn{করা} (sôndeh kôrā)}{সন্দেহ করা (sôndeh kôrā)}]{\bn{সন্দেহ} \bn{করা} (sôndeh kôrā) — to doubt}

\label{lesson:BN-C150-sondeh-kora}



\begin{quote}
Before the new one: say the Bengali for to hurry, then the Bengali for to delay.
\end{quote}

\subsection*{You'll want to know: \bn{সন্দেহ} \bn{করা}}
\textbf{\bn{সন্দেহ} \bn{করা}} — \emph{sôndeh kôrā} — \textquotedblleft{}to doubt\textquotedblright{}.

\begin{grammarlens}[title={What you've built}]
One more word for this chapter.
\end{grammarlens}

\subsection*{Your turn}
\begin{itemize}
\raggedright
  \item \emph{Say it:} \emph{sôndeh kôrā}
  \item \emph{Say it:} \emph{sôndeh kôrā}, once more
  \item \emph{Say it:} say \emph{sôndeh kôrā}
  \item \emph{Recall:} say the Bengali for to hurry, then the Bengali for to delay, then say \emph{sôndeh kôrā} again
\end{itemize}

\begin{tcolorbox}[breakable,colback=teal!4,colframe=teal!35!black,title={Before you move on}]
What does \bn{সন্দেহ} \bn{করা} mean? (\textquotedblleft{}To doubt\textquotedblright{}.) Say it once more.
\end{tcolorbox}
\section[\texorpdfstring{\bn{তর্ক} \bn{করা} (tôrkô kôrā)}{তর্ক করা (tôrkô kôrā)}]{\bn{তর্ক} \bn{করা} (tôrkô kôrā) — to argue}

\label{lesson:BN-C150-torko-kora}



\begin{quote}
Before the new one: say the Bengali for to delay, then the Bengali for to doubt.
\end{quote}

\subsection*{You'll want to know: \bn{তর্ক} \bn{করা}}
\textbf{\bn{তর্ক} \bn{করা}} — \emph{tôrkô kôrā} — \textquotedblleft{}to argue\textquotedblright{}.

\begin{grammarlens}[title={What you've built}]
One more word for this chapter.
\end{grammarlens}

\subsection*{Your turn}
\begin{itemize}
\raggedright
  \item \emph{Say it:} \emph{tôrkô kôrā}
  \item \emph{Say it:} \emph{tôrkô kôrā}, once more
  \item \emph{Say it:} say \emph{tôrkô kôrā}
  \item \emph{Recall:} say the Bengali for to delay, then the Bengali for to doubt, then say \emph{tôrkô kôrā} again
\end{itemize}

\begin{tcolorbox}[breakable,colback=teal!4,colframe=teal!35!black,title={Before you move on}]
What does \bn{তর্ক} \bn{করা} mean? (\textquotedblleft{}To argue\textquotedblright{}.) Say it once more.
\end{tcolorbox}
\section[\texorpdfstring{\bn{চালিয়ে} \bn{যাওয়া} (chāliye jāoyā)}{চালিয়ে যাওয়া (chāliye jāoyā)}]{\bn{চালিয়ে} \bn{যাওয়া} (chāliye jāoyā) — to carry on}

\label{lesson:BN-C150-chaliye-jaoya}



\begin{quote}
Before the new one: say the Bengali for to doubt, then the Bengali for to argue.
\end{quote}

\subsection*{You'll want to know: \bn{চালিয়ে} \bn{যাওয়া}}
\textbf{\bn{চালিয়ে} \bn{যাওয়া}} — \emph{chāliye jāoyā} — \textquotedblleft{}to carry on\textquotedblright{}.

\begin{grammarlens}[title={What you've built}]
One more word for this chapter.
\end{grammarlens}

\subsection*{Your turn}
\begin{itemize}
\raggedright
  \item \emph{Say it:} \emph{chāliye jāoyā}
  \item \emph{Say it:} \emph{chāliye jāoyā}, once more
  \item \emph{Say it:} say \emph{chāliye jāoyā}
  \item \emph{Recall:} say the Bengali for to doubt, then the Bengali for to argue, then say \emph{chāliye jāoyā} again
\end{itemize}

\begin{tcolorbox}[breakable,colback=teal!4,colframe=teal!35!black,title={Before you move on}]
What does \bn{চালিয়ে} \bn{যাওয়া} mean? (\textquotedblleft{}To carry on\textquotedblright{}.) Say it once more.
\end{tcolorbox}
\section[\texorpdfstring{\bn{হাঁচা} (hā̃chā)}{হাঁচা (hā̃chā)}]{\bn{হাঁচা} (hā̃chā) — to sneeze}

\label{lesson:BN-C150-hacha}



\begin{quote}
Before the new one: say the Bengali for to argue, then the Bengali for to carry on.
\end{quote}

\subsection*{You'll want to know: \bn{হাঁচা}}
\textbf{\bn{হাঁচা}} — \emph{hā̃chā} — \textquotedblleft{}to sneeze\textquotedblright{}.

\begin{grammarlens}[title={What you've built}]
One more word for this chapter.
\end{grammarlens}

\subsection*{Your turn}
\begin{itemize}
\raggedright
  \item \emph{Say it:} \emph{hā̃chā}
  \item \emph{Say it:} \emph{hā̃chā}, once more
  \item \emph{Say it:} say \emph{hā̃chā}
  \item \emph{Recall:} say the Bengali for to argue, then the Bengali for to carry on, then say \emph{hā̃chā} again
\end{itemize}

\begin{tcolorbox}[breakable,colback=teal!4,colframe=teal!35!black,title={Before you move on}]
What does \bn{হাঁচা} mean? (\textquotedblleft{}To sneeze\textquotedblright{}.) Say it once more.
\end{tcolorbox}
\section[\texorpdfstring{\bn{কাশা} (kāshā)}{কাশা (kāshā)}]{\bn{কাশা} (kāshā) — to cough}

\label{lesson:BN-C150-kasha}



\begin{quote}
Before the new one: say the Bengali for to carry on, then the Bengali for to sneeze.
\end{quote}

\subsection*{You'll want to know: \bn{কাশা}}
\textbf{\bn{কাশা}} — \emph{kāshā} — \textquotedblleft{}to cough\textquotedblright{}.

\begin{grammarlens}[title={What you've built}]
One more word for this chapter.
\end{grammarlens}

\subsection*{Your turn}
\begin{itemize}
\raggedright
  \item \emph{Say it:} \emph{kāshā}
  \item \emph{Say it:} \emph{kāshā}, once more
  \item \emph{Say it:} say \emph{kāshā}
  \item \emph{Recall:} say the Bengali for to carry on, then the Bengali for to sneeze, then say \emph{kāshā} again
\end{itemize}

\begin{tcolorbox}[breakable,colback=teal!4,colframe=teal!35!black,title={Before you move on}]
What does \bn{কাশা} mean? (\textquotedblleft{}To cough\textquotedblright{}.) Say it once more.
\end{tcolorbox}
